\section*{Exercise 3}

\textit{Task: Implementing the iterative deepening algorithm.}

\noindent Our task was to integrate the iterative deepening search algorithm into our game client in order to handle the time limits given by the server. The principle of iterative deepening is the continuous increase of the depth limit with each loop iteration. Running the alpha-beta pruning algorithm with the increased depth limit allows you to find an even better move until the time limit is reached.\\
We implemented the new \texttt{IterativeDeepening} class which includes the static \texttt{findMove(...)} method. With the help of a second class called \texttt{SearchInfo}, we keep track of the important information that should be available in the recursive function calls of alpha beta pruning. The most important attributes for this strategy are the given time limit and the timestamp from the start of the move computation. In every recursive call (for every node) we check whether the time has been reached. If so, we abort the evaluation and return the last computed move. We use a boolean value to achieve this timeout which is similar to an interrupt.However, with only one thread and recursive functions in Java, an ISR is almost impossible to implement. Therefore, we have also added an offset to have enough time left for all recursive calls to exit.\\
Before we came up with the abort mechanism, we tried to estimate the time required for the next iteration using the current search depth. This was also the main motivation behind the \texttt{SearchInfo} class. We started by computing various parameters to analyze the search tree's branching factor — such as the maximum, minimum, average, and average at a specific depth. We then combined these values and conducted several experiments in order to achieve gameplay without disqualifications due to timeouts.\\
However, we realized that finding an estimation method that works reliably in every situation is very difficult. While we were able to find some reasonably good estimations, disqualifications still occurred from time to time. The problem with our special version of Reversi is that the number of possible moves can change very rapidly depending on the current game state.\\
As a result, we switched to a safer approach using the abort mechanism. It might still be useful to combine this with an estimation in the future, in order to reduce wasted computation time and improve our chances of winning. But for now, we prefer to stay on the safe side.